\documentclass[11pt,a4paper]{article}
\usepackage[utf8]{inputenc}
\usepackage[margin=0.85in]{geometry}
\usepackage{amsmath,amssymb,amsthm}
\usepackage{xcolor}
\usepackage{tcolorbox}
\usepackage{tikz}
\usetikzlibrary{shapes,arrows,positioning,calc,decorations.pathreplacing}
\usepackage{booktabs}
\usepackage{array}
\usepackage{fancyhdr}
\usepackage{hyperref}
\usepackage{microtype}

% Color Palette Definitions
\definecolor{primarynavy}{RGB}{26, 54, 93}       % #1A365D - Deep Navy
\definecolor{secondaryteal}{RGB}{13, 148, 136}   % #0D9488 - Teal Accent
\definecolor{darkslate}{RGB}{45, 55, 72}        % #2D3748 - Body Text
\definecolor{lightbg}{RGB}{247, 250, 252}        % #F7FAFC - Soft Background
\definecolor{accentblue}{RGB}{49, 130, 206}      % #3182CE - Bright Blue
\definecolor{warmgold}{RGB}{214, 158, 46}        % #D69E2E - Gold Highlight

\hypersetup{
    colorlinks=true,
    linkcolor=primarynavy,
    citecolor=secondaryteal,
    urlcolor=accentblue,
    pdftitle={Lecture Notes: Multivariable Calculus \& Analysis},
    pdfauthor={Class Notes}
}

% Page Setup
\pagestyle{fancy}
\fancyhf{}
\fancyhead[L]{\small \textbf{\color{primarynavy}Multivariable Calculus \& Real Analysis}}
\fancyhead[R]{\small \textbf{\color{secondaryteal}MATMJ53}}
\fancyfoot[L]{\small \textcolor{primarynavy}{\href{https://www.sciqb.com}{\textbf{www.sciqb.com}}}}
\fancyfoot[C]{\thepage}
\fancyfoot[R]{\small \textcolor{secondaryteal}{\href{https://www.sciqb.com}{\textbf{SCIQB Platform}}}}
\renewcommand{\headrulewidth}{0.6pt}
\renewcommand{\footrulewidth}{0.4pt}

% Custom TColorBox Environments
\newtcolorbox{defbox}[1]{
    colback=lightbg,
    colframe=primarynavy,
    fonttitle=\bfseries\color{white},
    title={#1},
    boxrule=1pt,
    top=3mm, bottom=3mm, left=4mm, right=4mm
}

\newtcolorbox{formulabox}[1]{
    colback=secondaryteal!8!white,
    colframe=secondaryteal,
    fonttitle=\bfseries\color{white},
    title={#1},
    boxrule=1pt,
    top=3mm, bottom=3mm, left=4mm, right=4mm
}

\newtcolorbox{exbox}[1]{
    colback=white,
    colframe=accentblue,
    fonttitle=\bfseries\color{white},
    title={#1},
    boxrule=1pt,
    top=3mm, bottom=3mm, left=4mm, right=4mm
}

\setlength{\parindent}{0pt}
\setlength{\parskip}{6pt}

\begin{document}

% Title Banner
\begin{center}
    \tcbset{colframe=primarynavy,colback=primarynavy!5!white,arc=0mm}
    \begin{tcolorbox}[sharp corners, boxrule=1.5pt]
        \begin{center}
            {\Large \bfseries \color{primarynavy} LECTURE NOTES ON MULTIVARIABLE CALCULUS \& ANALYSIS}\\[4pt]
            {\large \bfseries \color{secondaryteal} Module 1: Functions of Several Variables \& Partial Differentiation \quad|\quad Module 2: Differentiability \& Taylor's Theorem}\\[6pt]
            {\small \textbf{Course Code:} MATMJ53 \quad|\quad \textbf{Platform:} \href{https://www.sciqb.com}{\color{accentblue}\textbf{www.sciqb.com}}}
        \end{center}
    \end{tcolorbox}
\end{center}

\vspace{0.3cm}

\tableofcontents
\newpage

\section{Module 1: Functions of Two Variables and Partial Differentiation}
\textit{Lecture Date: August 06, 2026}

Multivariable calculus extends single-variable calculus concepts to functions defined on $\mathbb{R}^n$.

\subsection{Limits and Continuity}

\begin{defbox}{Definition: Limit of a Function of Two Variables}
Let $f: D \subset \mathbb{R}^2 \to \mathbb{R}$ be a function. We say $\lim_{(x,y) \to (a,b)} f(x,y) = L$ if for every $\varepsilon > 0$, there exists $\delta > 0$ such that:
\begin{equation}
    0 < \sqrt{(x-a)^2 + (y-b)^2} < \delta \implies |f(x,y) - L| < \varepsilon
\end{equation}
\end{defbox}

\subsection{Partial Derivatives}

\begin{defbox}{Definitions: Partial Derivatives}
Let $f(x,y)$ be defined in a neighborhood of $(a,b)$.
\begin{itemize}
    \item \textbf{Partial Derivative with respect to $x$ at $(a,b)$:}
    \begin{equation}
        f_x(a,b) = \left. \frac{\partial f}{\partial x} \right|_{(a,b)} = \lim_{h \to 0} \frac{f(a+h, b) - f(a,b)}{h}
    \end{equation}
    
    \item \textbf{Partial Derivative with respect to $y$ at $(a,b)$:}
    \begin{equation}
        f_y(a,b) = \left. \frac{\partial f}{\partial y} \right|_{(a,b)} = \lim_{k \to 0} \frac{f(a, b+k) - f(a,b)}{k}
    \end{equation}
\end{itemize}
\end{defbox}

\section{Module 2: Differentiability and Total Differentials}

\subsection{Definition of Differentiability}

\begin{formulabox}{Differentiability of $f(x,y)$ at $(a,b)$}
A function $f(x,y)$ is said to be \textbf{differentiable} at $(a,b)$ if the increment $\Delta f = f(a+h, b+k) - f(a,b)$ can be expressed as:
\begin{equation}
    f(a+h, b+k) - f(a,b) = h f_x(a,b) + k f_y(a,b) + h \phi(h,k) + k \psi(h,k)
\end{equation}
where $\phi(h,k) \to 0$ and $\psi(h,k) \to 0$ as $(h,k) \to (0,0)$.
\end{formulabox}

\begin{exbox}{Theorem: Differentiability Implies Continuity}
\textbf{Theorem:} If a function $f(x,y)$ is differentiable at $(a,b)$, then $f(x,y)$ is continuous at $(a,b)$.

\textbf{Proof:} Taking the limit as $(h,k) \to (0,0)$ in the differentiability increment formula yields:
$$\lim_{(h,k) \to (0,0)} f(a+h, b+k) = f(a,b)$$
Thus $f$ is continuous at $(a,b)$. \qed
\end{exbox}

\subsection{Taylor's Theorem for Functions of Two Variables}

\begin{defbox}{Taylor Series Expansion}
If $f(x,y)$ and its partial derivatives up to order $n$ are continuous in a neighborhood of $(a,b)$, then:
\begin{equation}
    f(a+h, b+k) = f(a,b) + \left( h \frac{\partial}{\partial x} + k \frac{\partial}{\partial y} \right) f(a,b) + \frac{1}{2!} \left( h \frac{\partial}{\partial x} + k \frac{\partial}{\partial y} \right)^2 f(a,b) + \dots + R_n
\end{equation}
\end{defbox}

\end{document}
